\documentclass[a4paper,11pt]{article}
\usepackage[english]{babel}
\usepackage[utf8]{inputenc}
\usepackage[T1]{fontenc}
\usepackage[a4paper,margin=1.5cm,top=1.2cm,bottom=1.2cm]{geometry}
\usepackage{graphicx}
\usepackage{amsmath}
\usepackage{booktabs}
\usepackage{siunitx}
\usepackage{float}
\usepackage{microtype}
\usepackage{newtxtext,newtxmath}
\usepackage{setspace}
\usepackage{fancyhdr}
\setstretch{1.15}

\bibliographystyle{apsrev}
\setcounter{secnumdepth}{3}
\pagestyle{fancy}
\fancyhf{}
\fancyfoot[C]{\thepage}
\renewcommand{\headrulewidth}{0pt}
\renewcommand{\footrulewidth}{0pt}
\setlength{\headheight}{14pt}
\setlength{\footskip}{14pt}

\begin{document}

\title{Determinacion de la longitud de onda mediante un biprisma de Fresnel}
\author{Practicas de Optica}
\maketitle

\begin{abstract}
Se determina experimentalmente la longitud de onda de la luz mediante un biprisma
de Fresnel. El analisis utiliza las distancias medidas entre las fuentes virtuales
y la pantalla, junto con la separacion entre franjas. Los resultados se comparan
con el valor teorico de $589\,\mathrm{nm}$.
\end{abstract}
COMENTARIO IMPORTANTE: YO LE HE LLAMADO PRIMERA PARTE PERO DIRIA DE CATEGORIZARLO EN 3 MEDIDAS VS 1 MEDIDA, ADEMAS AL FINAL COMO EN ELC ASO D E1 SE REDONDEA POR EXCEOS Y POR DEFECTO TAMBIEN SALE UN RESULTADO CORRECOT PERO PODEMOS VER QUE 
LA INCERTIDUMBRE ES MAYOR.
\section{Introduccion y objetivos}

\subsection{Fundamento teorico}
Las rendijas de luz con las que trabajamos son la imagen virtual de un biprisma de fresnel que iluminamos mediante una lamaprado de soido.
Como es una fuente de luz monocromatica, sus dos vectores eléctricos son paraleleos, al ser haces de luz coherentes su amplitud y frecuencia son iguales
La expresion del campo electrico es
\begin{equation}
E = E_0 e^{i(\omega t - k x + \phi)}
\end{equation}
En este patron de interferencia se tiene que la distancia entre dos maximos consecutivos,
es decir, dos franjas visibles, es
\begin{equation}
x_m-x_{m-1}=\lambda\frac{D}{d}.
\end{equation}
Donde $x_m$ es la posicion del maximo $m$, $\lambda$ la longitud de onda del sodio,
$D$ la distancia entre el punto medio de las fuentes puntuales y el plano $xy$ que
contiene el punto de observacion, y $d$ la separacion de las fuentes puntuales.

El problema de la ecuacion es que $D$ no se puede medir en el banco optico, ya que
requeriria conocer con exactitud la posicion de la imagen virtual producida por el
biprisma. Por lo que lo determinaremos a partir de otras medidas, para lo cual será necesario obtener $d$. 
Para ello, usaremos que la ecuacion de Gauss nos permite obtener dos
posiciones distintas de la lente para las cuales se forma una imagen nitida sobre
la pantalla. Asi que, tras determinar estas posiciones, podemos medir la distancia
de las franjas con la lente en estas posiciones.

Podemos ver que los aumentos en estas posiciones serian $d_1/d$ y $d_2/d$, y ademas
la ecuacion de Gauss nos dice que un aumento resulta el inverso del otro, ya que la
distancia objeto de la posicion 1 es la distancia imagen de la posicion 2 y viceversa.
Se obtiene entonces que $d_1/d=d/d_2$, de donde se obtiene la expresion para $d$:
\begin{equation}
d=\sqrt{d_1d_2}.
\end{equation}
Esto proporciona una forma de obtener este parametro, que usaremos para obtener $D$
con la siguiente ecuacion:
\begin{equation}
D=f'\left(2-\frac{d^2+d_1^2}{d d_1}\right),
\end{equation}
donde $d_1$ es la menor de las distancias medidas anteriormente.


\subsection{Metodologia de medida}
\subsubsection{Montaje experimental}
Experimentalmente, entonces, tenemos nuestra fuente de luz monocromatica, que en
nuestro caso era una lampara de sodio, el biprisma de Fresnel y un ocular
micrometrico que nos permitia observar la interferencia.
Cabe recalcar que en la practica se presentaron dos incovenientes importantes:
El viñeteo de la lampara de sodio era considerable. Esto nos llevo a poner una
cartulina negra, de tal forma que el area de observacion fuese mas clara.

Fuera de la zona central de interferencia, estas no estaban equiespaciadas, lo cual
es una contradiccion de nuestro modelo. Probablemente se deba a imperfecciones en
el material, como pueden ser anisotropias por burbujas locales o danos estructurales.
De tal forma que, para subsanar este problema, decidimos tomar siempre las medidas
en torno a las franjas centrales, que si estaban equiespaciadas entre ellas.

En la zona de interferencia se observan las franjas equiespaciadas en el centro y
las no equiespaciadas en los extremos, aunque en nuestro caso este problema era mas
pronunciado.

\subsubsection{Mediciones}
Una vez calibrado el dispositivo experimental y subsanados los errores, se procedió a tomar las medidas de la distancia entre las franjas.
Este proceso fue llevado a cabo de dos formas distintas: en la primera se tomaron
tres medidas de la distancia entre tres franjas, y en la segunda tomamos una franja
para dos lentes con distancias focales de 15 y 20 cm, respectivamente.

La meta de este proceso era obtener la longitud de onda de la luz, asi como
cuantificar el error que se comete al medir diferentes cantidades de franjas.
Pese a que en un primer vistazo pudiera parecer que el hecho de medir tres franjas
es mas optimo solamente por la menor incertidumbre intrinseca, presenta una ventaja
adicional: corrige los errores de medicion del cientifico, ya que muchas veces
discernir donde acaba y termina una sola franja es complicado y se comete un error.
Al hacerlo sobre tres franjas, este error representa un porcentaje mucho menor del total.

En retrospectiva, quizas el metodo mas eficiente sea contar en primera instancia cual
es el numero maximo de franjas equidistantes y hacer la medida sobre este numero para
reducir la incertidumbre.

Cabe destacar que en la segunda parte, con la lente de 15 cm, realizamos las medidas
en el propio ocular, prescindiendo del tornillo micrometrico, lo cual nos indujo a
un error adicional.

Para obtener la longitud de onda se utiliza
\begin{equation}
s=3\lambda\frac{D}{d}.
\end{equation}
La incertidumbre que acompana a cada valor de $\lambda$ es la incertidumbre expandida,
calculada considerando conjuntamente el error de tipo A y el error de tipo B.







\section{Resultados y analisis de datos}

\subsection{Resultados de la longitud de onda}
La Tabla~\ref{tab:lambda} resume los resultados obtenidos para cada medida.

\begin{table}[H]
\centering
\begin{tabular}{lccc}
\toprule
\textbf{Parte} & $\lambda_1$ & $\lambda_2$ & $\lambda_3$ \\
\midrule
Parte 1 & $577\pm87$ & $594\pm95$ & $569\pm92$ \\
Parte 2 & $693\pm33$ & $808\pm33$ & $750\pm33$ \\
Parte 3 & $491\pm130$ & $508\pm122$ & $505\pm146$ \\
\bottomrule
\end{tabular}
\caption{Longitudes de onda individuales en nm.}
\label{tab:lambda}
\end{table}

\subsection{Promedios}
\begin{table}[H]
\centering
\begin{tabular}{lc}
\toprule
\textbf{Parte} & \textbf{Promedio de $\lambda$ (nm)} \\
\midrule
Parte 1 & $580\pm53$ \\
Partes 2 y 3 juntas & $626\pm70$ \\
\bottomrule
\end{tabular}
\caption{Promedios calculados combinando la dispersion entre medidas y sus incertidumbres propagadas.}
\label{tab:promedios}
\end{table}

Juntando las tres medidas de la parte 2 y las tres medidas de la parte 3, la media
de las seis medidas de la longitud de onda es
\begin{equation}
\overline{\lambda}_{2+3}
 = \frac{693+808+750+491+508+505}{6}\,\mathrm{nm}
 = (626\pm70)\,\mathrm{nm}.
\end{equation}

El resultado de la primera parte, $580\pm53\,\mathrm{nm}$, es compatible con
este valor dentro de la incertidumbre experimental.

El valor teorico de referencia es $\lambda_\mathrm{teo}=589\,\mathrm{nm}$.

\section{Discusion y conclusiones}
La primera parte presenta el mejor acuerdo con el valor teorico de la fuente de
sodio. La media conjunta de las partes 2 y 3 es $(626\pm70)\,\mathrm{nm}$.
Las diferencias entre las partes pueden deberse a la sensibilidad de $D$ frente a
las distancias medidas y a la incertidumbre relativamente grande de la mirilla
usada en la tercera parte. La propagacion incluye las incertidumbres de $s$,
$d'_1$ y $d'_2$ en $d$, $D$ y $\lambda$.

\begin{equation}
\lambda_1 = (580\pm53)\,\mathrm{nm}, \qquad
\overline{\lambda}_{2+3} = (626\pm70)\,\mathrm{nm}.
\end{equation}

\end{document}
